\chapter{Resultados de simulación}
\label{chap:resultados}

\section{Intervalo de operación del actuador}
\label{sec:intervalo}

El voltaje que puede aplicar el actuador está acotado a $u\in[0;\,3{,}5]$~V, y esa cota determina en qué posiciones puede operar el sistema. Para sostener el imán en $x_1$ se necesita el voltaje de equilibrio $u_e(x_1)=mgb(a-x_1)^4$, que crece rápidamente al alejarse de la bobina: $1{,}58$~V en $-2$~cm, $2{,}39$~V en $-3$~cm, $3{,}48$~V en $-4$~cm y $4{,}91$~V en $-5$~cm. Para mover el imán hay que vencer además la fricción estática, lo que exige llegar a $u_e(1\pm F_s/mg)$, es decir, a un $14{,}5$\,\% por encima o por debajo del equilibrio (sección~\ref{sec:analisis_ciclo}).

En $x_1=-4$~cm, el punto en que se linealizó el modelo para el diseño, el imán puede sostenerse con $3{,}48$~V, pero no desplazarse: para despegarlo hacia la bobina se requieren unos $3{,}99$~V. Para delimitar el intervalo útil se simularon 7503 escalones con distintas posiciones iniciales y finales, con el PII y con el PI, en presencia de atascamiento-deslizamiento. La tabla~\ref{tab:intervalo} resume el resultado con el PII.

\begin{table}[ht]
\centering
\caption{Referencias alcanzables con $u\le3{,}5$~V desde distintas posiciones iniciales (PII, con atascamiento-deslizamiento).}
\label{tab:intervalo}
\begin{tabular}{|c|c|}
\hline
Posición inicial [cm] & Referencias alcanzables [cm] \\ \hline
$-4{,}0$ & ninguna: el imán se sostiene pero no se mueve \\ \hline
$-3{,}5$ & de $-3{,}6$ a $-2{,}0$ \\ \hline
$-3{,}0$ & de $-3{,}5$ a $-1{,}5$ \\ \hline
$-2{,}0$ & de $-3{,}3$ a $-0{,}5$ \\ \hline
$-1{,}0$ & de $-2{,}5$ a $+0{,}5$ \\ \hline
\end{tabular}
\end{table}

Los cambios de referencia son viables a partir de unos $-3{,}6$~cm. Por eso, las pruebas de este capítulo se realizan en el intervalo de $-3$ a $-2$~cm, con un escalón de $-2$ a $-3$~cm que tiene margen respecto a ambos límites. El controlador diseñado en $-4$~cm sigue siendo estable en ese intervalo, con márgenes de fase de $54{,}5^\circ$ en $-3$~cm, $51{,}1^\circ$ en $-2{,}5$~cm y $47{,}4^\circ$ en $-2$~cm.

\section{Simulación no lineal sin atascamiento-deslizamiento}

Antes de incorporar el efecto atascamiento-deslizamiento, el controlador PII se evaluó sobre el modelo no lineal del levitador considerando únicamente la fuerza electromagnética no lineal y la fricción viscosa,
\begin{equation}
m\ddot{x} = \frac{u}{b\,(a-x)^4} - m g - m c_1 \dot{x} ,
\label{eq:nl_sin_ad}
\end{equation}
es decir, sin fricción estática (modelo de Karnopp) ni zona muerta en el actuador. Se emplearon las mismas condiciones que en la prueba de regulación de la sección siguiente: escalón de referencia de $-2$ a $-3$~cm en $t=10$~s, señal de control acotada a $u\in[0;\,3{,}5]$~V y controlador inicializado en el estado estacionario correspondiente a $x_1=-2$~cm. Para facilitar la comparación, en cada figura se superpone el resultado con atascamiento-deslizamiento.

\begin{figure}[ht]
    \centering
    % Estilo común de las figuras sin/con atascamiento-deslizamiento (pgfplots).
\pgfplotsset{
  sinad/.style={
    width=0.92\textwidth, height=5.2cm, xmin=0, xmax=40,
    xlabel={tiempo [\si{s}]},
    grid=major, grid style={gray!25},
    tick label style={font=\small}, label style={font=\small}, legend style={font=\small},
    /pgf/number format/use comma,
    scaled ticks=false,
    y tick label style={/pgf/number format/fixed},
    table/col sep=comma,
    every axis plot/.append style={line join=round},
  },
  detalle/.style={height=4cm, axis background/.style={fill=ctl2!4}},
  sinAD/.style={tray, line width=1.1pt},
  conAD/.style={ctl2, line width=0.7pt},
}
\definecolor{tray}{RGB}{31,94,166}
\definecolor{refc}{RGB}{40,40,40}
\definecolor{ctl2}{RGB}{196,96,40}

\begin{tikzpicture}
\begin{axis}[sinad, name=principal, xlabel={}, ylabel={posición [\si{cm}]}, ymin=-3.35, ymax=-1.85,
             legend style={at={(0.99,0.97)}, anchor=north east}, legend cell align=left]
  \addplot[conAD] table[x=t, y=x_ad] {images/sin_ad_results/tikz/comparacion.csv};
  \addlegendentry{con AD}
  \addplot[sinAD] table[x=t, y=x] {images/sin_ad_results/tikz/comparacion.csv};
  \addlegendentry{sin AD}
  \addplot[refc, densely dashed, line width=0.8pt] table[x=t, y=r] {images/sin_ad_results/tikz/comparacion.csv};
  \addlegendentry{referencia}
  \fill[ctl2, opacity=0.15] (axis cs:9.8,-3.35) rectangle (axis cs:12,-1.85);
\end{axis}
\begin{axis}[sinad, detalle, at={(principal.south west)}, anchor=north west, yshift=-0.45cm,
             xmin=9.8, xmax=12, ymin=-3.3, ymax=-1.9, ylabel={detalle [\si{cm}]}]
  \addplot[conAD, line width=1.6pt] table[x=t, y=x_ad] {images/sin_ad_results/tikz/detalle.csv};
  \addplot[sinAD, line width=0.8pt] table[x=t, y=x] {images/sin_ad_results/tikz/detalle.csv};
  \addplot[refc, densely dashed, line width=0.8pt] table[x=t, y=r] {images/sin_ad_results/tikz/detalle.csv};
\end{axis}
\end{tikzpicture}

    \caption[Sin AD: posición]{Posición del levitante con el modelo no lineal sin atascamiento-deslizamiento (sin AD) y con él (con AD). Abajo, detalle del transitorio entre 9,8 y 12~s.}
    \label{fig:sinad_pos}
\end{figure}

\begin{figure}[ht]
    \centering
    % Estilo común de las figuras sin/con atascamiento-deslizamiento (pgfplots).
\pgfplotsset{
  sinad/.style={
    width=0.92\textwidth, height=5.2cm, xmin=0, xmax=40,
    xlabel={tiempo [\si{s}]},
    grid=major, grid style={gray!25},
    tick label style={font=\small}, label style={font=\small}, legend style={font=\small},
    /pgf/number format/use comma,
    scaled ticks=false,
    y tick label style={/pgf/number format/fixed},
    table/col sep=comma,
    every axis plot/.append style={line join=round},
  },
  detalle/.style={height=4cm, axis background/.style={fill=ctl2!4}},
  sinAD/.style={tray, line width=1.1pt},
  conAD/.style={ctl2, line width=0.7pt},
}
\definecolor{tray}{RGB}{31,94,166}
\definecolor{refc}{RGB}{40,40,40}
\definecolor{ctl2}{RGB}{196,96,40}

\begin{tikzpicture}
\begin{axis}[sinad, ylabel={velocidad [\si{cm/s}]},
             legend style={at={(0.99,0.03)}, anchor=south east}, legend cell align=left]
  \addplot[conAD] table[x=t, y=v_ad] {images/sin_ad_results/tikz/comparacion.csv};
  \addlegendentry{con AD}
  \addplot[sinAD] table[x=t, y=v] {images/sin_ad_results/tikz/comparacion.csv};
  \addlegendentry{sin AD}
\end{axis}
\end{tikzpicture}

    \caption[Sin AD: velocidad]{Velocidad del levitante sin y con atascamiento-deslizamiento.}
    \label{fig:sinad_vel}
\end{figure}

\begin{figure}[ht]
    \centering
    % Estilo común de las figuras sin/con atascamiento-deslizamiento (pgfplots).
\pgfplotsset{
  sinad/.style={
    width=0.92\textwidth, height=5.2cm, xmin=0, xmax=40,
    xlabel={tiempo [\si{s}]},
    grid=major, grid style={gray!25},
    tick label style={font=\small}, label style={font=\small}, legend style={font=\small},
    /pgf/number format/use comma,
    scaled ticks=false,
    y tick label style={/pgf/number format/fixed},
    table/col sep=comma,
    every axis plot/.append style={line join=round},
  },
  detalle/.style={height=4cm, axis background/.style={fill=ctl2!4}},
  sinAD/.style={tray, line width=1.1pt},
  conAD/.style={ctl2, line width=0.7pt},
}
\definecolor{tray}{RGB}{31,94,166}
\definecolor{refc}{RGB}{40,40,40}
\definecolor{ctl2}{RGB}{196,96,40}

\begin{tikzpicture}
\begin{axis}[sinad, name=principal, xlabel={}, ylabel={error [\si{cm}]},
             legend style={at={(0.99,0.03)}, anchor=south east}, legend cell align=left]
  \addplot[conAD] table[x=t, y=e_ad] {images/sin_ad_results/tikz/comparacion.csv};
  \addlegendentry{con AD}
  \addplot[sinAD] table[x=t, y=e] {images/sin_ad_results/tikz/comparacion.csv};
  \addlegendentry{sin AD}
  \fill[ctl2, opacity=0.15] ({axis cs:15,0} |- {rel axis cs:0,0}) rectangle ({axis cs:40,0} |- {rel axis cs:0,1});
\end{axis}
\begin{axis}[sinad, detalle, at={(principal.south west)}, anchor=north west, yshift=-0.45cm,
             xmin=15, xmax=40, ymin=-0.03, ymax=0.03, ylabel={detalle [\si{cm}]}]
  \addplot[refc!60, thin] coordinates {(15,0) (40,0)};
  \addplot[conAD] table[x=t, y=e_ad] {images/sin_ad_results/tikz/comparacion.csv};
  \addplot[sinAD] table[x=t, y=e] {images/sin_ad_results/tikz/comparacion.csv};
\end{axis}
\end{tikzpicture}

    \caption[Sin AD: error]{Error de posición sin y con atascamiento-deslizamiento. Abajo, detalle entre 15 y 40~s.}
    \label{fig:sinad_error}
\end{figure}

\begin{figure}[ht]
    \centering
    % Estilo común de las figuras sin/con atascamiento-deslizamiento (pgfplots).
\pgfplotsset{
  sinad/.style={
    width=0.92\textwidth, height=5.2cm, xmin=0, xmax=40,
    xlabel={tiempo [\si{s}]},
    grid=major, grid style={gray!25},
    tick label style={font=\small}, label style={font=\small}, legend style={font=\small},
    /pgf/number format/use comma,
    scaled ticks=false,
    y tick label style={/pgf/number format/fixed},
    table/col sep=comma,
    every axis plot/.append style={line join=round},
  },
  detalle/.style={height=4cm, axis background/.style={fill=ctl2!4}},
  sinAD/.style={tray, line width=1.1pt},
  conAD/.style={ctl2, line width=0.7pt},
}
\definecolor{tray}{RGB}{31,94,166}
\definecolor{refc}{RGB}{40,40,40}
\definecolor{ctl2}{RGB}{196,96,40}

\begin{tikzpicture}
\begin{axis}[sinad, ylabel={señal de control [\si{V}]}, ymin=-0.2, ymax=3.8,
             legend style={at={(0.99,0.03)}, anchor=south east}, legend cell align=left, legend columns=3]
  \addplot[conAD] table[x=t, y=u_ad] {images/sin_ad_results/tikz/comparacion.csv};
  \addlegendentry{con AD}
  \addplot[sinAD] table[x=t, y=u] {images/sin_ad_results/tikz/comparacion.csv};
  \addlegendentry{sin AD}
  \addplot[red!70!black, densely dotted, line width=0.8pt] coordinates {(0,3.5) (40,3.5)};
  \addlegendentry{$u_{\max}$}
\end{axis}
\end{tikzpicture}

    \caption[Sin AD: señal de control]{Señal de control sin y con atascamiento-deslizamiento.}
    \label{fig:sinad_ctrl}
\end{figure}

La figura~\ref{fig:sinad_pos} muestra que, sin atascamiento-deslizamiento, el lazo cerrado no lineal alcanza la nueva referencia con un sobrepaso de $23{,}7$\,\% a los $0{,}15$~s del escalón y entra en la banda de $\pm 2$\,\% del cambio de referencia en $0{,}48$~s. Este sobrepaso se encuentra entre los que predice el modelo linealizado en $x_1^*=-3$~cm ($19{,}2$\,\%) y en $x_1^*=-2$~cm ($24{,}7$\,\%), lo que confirma la coherencia entre el modelo empleado para el diseño y la planta no lineal. La velocidad (figura~\ref{fig:sinad_vel}) se anula una vez concluido el transitorio, el error (figura~\ref{fig:sinad_error}) converge a cero y la señal de control (figura~\ref{fig:sinad_ctrl}) converge al voltaje de equilibrio $u_e=2{,}3934$~V correspondiente a $x_1=-3$~cm, sin alcanzar la saturación superior.

La comparación con el caso con atascamiento-deslizamiento muestra que el transitorio es prácticamente idéntico (sobrepaso de $23{,}6$\,\%): la diferencia aparece en estado estacionario, donde la fricción estática y la zona muerta sustituyen la convergencia asintótica por un ciclo límite de $0{,}033$~cm pico a pico, que se describe en la sección siguiente y se analiza en la sección~\ref{sec:analisis_ciclo}. Con ello se cumple el tercer objetivo específico.

\FloatBarrier
\newpage

\section{Regulación}

La prueba de regulación aplica un escalón de referencia de $-2$ a $-3$~cm en $t=10$~s al lazo con el PII, la zona muerta y la fricción de Karnopp, con la señal de control acotada a $u\in[0;\,3{,}5]$~V. Se simulan 40~s con un periodo de muestreo de 1~ms; el controlador parte del estado estacionario correspondiente a $x_1=-2$~cm, de modo que antes del escalón el levitante permanece en reposo.

La figura~\ref{posicion} muestra la respuesta. El levitante alcanza un mínimo de $-3{,}237$~cm a los $0{,}15$~s del escalón, lo que equivale a un sobrepaso de $23{,}6$\,\%, y entra en la banda de $\pm0{,}05$~cm alrededor de la referencia a los $0{,}40$~s. A partir de ahí no se detiene sobre la referencia: oscila en un ciclo límite de $0{,}033$~cm pico a pico, con un error máximo de $0{,}018$~cm y un error medio de $6\times10^{-5}$~cm en los últimos 10~s. La velocidad (figura~\ref{velocidad}) alterna pulsos de deslizamiento con intervalos en cero; el levitante permanece atascado el 74\,\% del tiempo en estado estacionario. El error (figura~\ref{error}) refleja el mismo comportamiento.

La señal de control (figura~\ref{senal_control}) cae momentáneamente a 0~V en el instante del escalón, que es el límite inferior del actuador, para dejar que el imán se aleje de la bobina; alcanza un máximo de $3{,}28$~V durante el transitorio y en estado estacionario oscila en diente de sierra alrededor de su valor medio de $2{,}41$~V, sin llegar al límite superior. El mecanismo del ciclo límite se analiza en la sección~\ref{sec:analisis_ciclo}.

\begin{figure}[ht]
    \centering
    % Estilo común de las figuras de regulación (pgfplots). Se carga con \input antes de cada figura.
\pgfplotsset{
  regulacion/.style={
    width=0.92\textwidth, height=5.6cm,
    xmin=0, xmax=40,
    xlabel={tiempo [\si{s}]},
    grid=major, grid style={gray!25},
    tick label style={font=\small}, label style={font=\small}, legend style={font=\small},
    /pgf/number format/use comma,
    scaled y ticks=false,
    y tick label style={/pgf/number format/fixed},
    table/col sep=comma,
    every axis plot/.append style={line join=round},
  },
}
\definecolor{tray}{RGB}{31,94,166}
\definecolor{refc}{RGB}{40,40,40}
\definecolor{ctl2}{RGB}{196,96,40}

\begin{tikzpicture}
\begin{axis}[regulacion, ylabel={posición [\si{cm}]}, ymin=-3.35, ymax=-1.85,
             legend pos=south west, legend cell align=left]
  \addplot[tray, line width=1pt] table[x=t, y=x] {images/regulation_results/tikz/posicion.csv};
  \addlegendentry{posición $x_1$}
  \addplot[refc, densely dashed, line width=0.8pt] table[x=t, y=r] {images/regulation_results/tikz/posicion.csv};
  \addlegendentry{referencia}
  \coordinate (zoomA) at (axis cs:9.8,-3.27);
  \coordinate (zoomB) at (axis cs:14,-2.95);
  \coordinate (inset) at (axis cs:39.2,-1.93);
\end{axis}
\draw[gray, thin] (zoomA) rectangle (zoomB);
\begin{axis}[regulacion, at={(inset)}, anchor=north east, width=0.5\textwidth, height=3.1cm,
             xmin=9.8, xmax=14, ymin=-3.27, ymax=-2.95, xlabel={}, ylabel={},
             tick label style={font=\scriptsize}, axis background/.style={fill=white},
             xtick={10,11,12,13,14}, ytick={-3.2,-3.1,-3}, name=zoom]
  \addplot[tray, line width=0.9pt] table[x=t, y=x] {images/regulation_results/tikz/posicion.csv};
  \addplot[refc, densely dashed, line width=0.7pt] table[x=t, y=r] {images/regulation_results/tikz/posicion.csv};
\end{axis}
\draw[gray, thin, densely dotted] (zoomB) -- (zoom.south west);
\end{tikzpicture}

    \caption[Posición (regulación)]{Posición del levitante ante un escalón de referencia de $-2$ a $-3$~cm en $t=10$~s (PII, $u\in[0;\,3{,}5]$~V). El recuadro amplía el transitorio y el ciclo límite de atascamiento-deslizamiento.}
    \label{posicion}
\end{figure}

\begin{figure}[ht]
    \centering
    % Estilo común de las figuras de regulación (pgfplots). Se carga con \input antes de cada figura.
\pgfplotsset{
  regulacion/.style={
    width=0.92\textwidth, height=5.6cm,
    xmin=0, xmax=40,
    xlabel={tiempo [\si{s}]},
    grid=major, grid style={gray!25},
    tick label style={font=\small}, label style={font=\small}, legend style={font=\small},
    /pgf/number format/use comma,
    scaled y ticks=false,
    y tick label style={/pgf/number format/fixed},
    table/col sep=comma,
    every axis plot/.append style={line join=round},
  },
}
\definecolor{tray}{RGB}{31,94,166}
\definecolor{refc}{RGB}{40,40,40}
\definecolor{ctl2}{RGB}{196,96,40}

\begin{tikzpicture}
\begin{axis}[regulacion, ylabel={velocidad [\si{cm/s}]}]
  \addplot[tray, line width=0.8pt] table[x=t, y=v] {images/regulation_results/tikz/velocidad.csv};
\end{axis}
\end{tikzpicture}

    \caption[Velocidad (regulación)]{Velocidad del levitante. Los tramos en cero corresponden al régimen de atascamiento.}
    \label{velocidad}
\end{figure}

\begin{figure}[ht]
    \centering
    % Estilo común de las figuras de regulación (pgfplots). Se carga con \input antes de cada figura.
\pgfplotsset{
  regulacion/.style={
    width=0.92\textwidth, height=5.6cm,
    xmin=0, xmax=40,
    xlabel={tiempo [\si{s}]},
    grid=major, grid style={gray!25},
    tick label style={font=\small}, label style={font=\small}, legend style={font=\small},
    /pgf/number format/use comma,
    scaled y ticks=false,
    y tick label style={/pgf/number format/fixed},
    table/col sep=comma,
    every axis plot/.append style={line join=round},
  },
}
\definecolor{tray}{RGB}{31,94,166}
\definecolor{refc}{RGB}{40,40,40}
\definecolor{ctl2}{RGB}{196,96,40}

\begin{tikzpicture}
\begin{axis}[regulacion, ylabel={error $e=r-x_1$ [\si{cm}]}]
  \addplot[refc!60, thin] coordinates {(0,0) (40,0)};
  \addplot[tray, line width=0.9pt] table[x=t, y=e] {images/regulation_results/tikz/error.csv};
\end{axis}
\end{tikzpicture}

    \caption[Error (regulación)]{Error de posición $e=r-x_1$.}
    \label{error}
\end{figure}

\begin{figure}[ht]
    \centering
    % Estilo común de las figuras de regulación (pgfplots). Se carga con \input antes de cada figura.
\pgfplotsset{
  regulacion/.style={
    width=0.92\textwidth, height=5.6cm,
    xmin=0, xmax=40,
    xlabel={tiempo [\si{s}]},
    grid=major, grid style={gray!25},
    tick label style={font=\small}, label style={font=\small}, legend style={font=\small},
    /pgf/number format/use comma,
    scaled y ticks=false,
    y tick label style={/pgf/number format/fixed},
    table/col sep=comma,
    every axis plot/.append style={line join=round},
  },
}
\definecolor{tray}{RGB}{31,94,166}
\definecolor{refc}{RGB}{40,40,40}
\definecolor{ctl2}{RGB}{196,96,40}

\begin{tikzpicture}
\begin{axis}[regulacion, ylabel={señal de control [\si{V}]}, ymin=-0.2, ymax=3.8,
             legend pos=south east, legend cell align=left]
  \addplot[ctl2!45, line width=1.6pt] table[x=t, y=u_antes] {images/regulation_results/tikz/control.csv};
  \addlegendentry{antes de la zona muerta}
  \addplot[tray, line width=0.7pt] table[x=t, y=u_desp] {images/regulation_results/tikz/control.csv};
  \addlegendentry{después de la zona muerta}
  \addplot[red!70!black, densely dotted, line width=0.8pt] coordinates {(0,3.5) (40,3.5)};
  \addlegendentry{$u_{\max}=\SI{3,5}{V}$}
\end{axis}
\end{tikzpicture}

    \caption[Señal de control (regulación)]{Señal de control antes y después de la zona muerta.}
    \label{senal_control}
\end{figure}

\FloatBarrier
\newpage

\section{Seguimiento}

Con el fin de evaluar el desempeño dinámico del controlador PII más allá del problema de regulación,
se realizaron experimentos de seguimiento en los que la referencia de posición varía en el tiempo.
Estas pruebas permiten analizar la interacción entre la acción de doble integración del controlador,
la no linealidad de zona muerta del actuador y el efecto atascamiento-deslizamiento modelado con
la fricción de Karnopp, todo ello en el intervalo de operación de $-3$ a $-2$ cm, en el que el actuador, limitado a $u\in[0;\,3{,}5]$~V, puede vencer la fricción estática en ambos sentidos.
El precedente experimental más cercano es el presentado por Alcántara et al.~\cite{alcantara},
quienes demostraron en tiempo real tanto regulación como seguimiento para el mismo dispositivo
ECP-730 en configuración repulsiva (estable), empleando un esquema de linealización por
retroalimentación de estados combinado con un lazo externo con acción integral y control conmutado.
El presente trabajo extiende ese resultado al caso de configuración atractiva (inestable),
donde los puntos de equilibrio son inherentemente inestables y la tarea de seguimiento resulta
más exigente al requerir estabilización simultánea.

%=============================================================================
\subsection{Seguimiento de referencia senoidal}
%=============================================================================

La referencia de posición utilizada en esta prueba tiene la forma

\begin{equation}
  r(t) = x_0 + A\,\sin(2\pi f\, t),
  \label{eq:ref_senoidal}
\end{equation}

con valor nominal $x_0 = -2{,}5$ cm, amplitud $A = 0{,}5$ cm y frecuencia $f = 0{,}001$ Hz
(periodo $T = 1000$ s).
La amplitud fue seleccionada de modo que la referencia excursiona simétricamente entre
$-3$ cm y $-2$ cm, el mismo intervalo de la prueba de regulación. Se simularon dos periodos (2000~s).

\paragraph{Comportamiento esperado según el modelo lineal.}
Sin atascamiento-deslizamiento, el lazo cerrado lineal haría que la posición siguiera la referencia senoidal con un desfase pequeño y una atenuación de amplitud mínima. Si el sistema linealizado en lazo cerrado tiene función de transferencia $T(s)$, con $T(j\omega_0)=|T|\,e^{j\varphi}$, el error en estado estacionario ante una referencia senoidal de frecuencia $\omega_0 = 2\pi f$ es

\begin{equation}
  e_{ss}(t) = A\bigl[1 - |T|\cos\varphi\bigr]\sin(\omega_0 t)
              - A\,|T|\sin\varphi\,\cos(\omega_0 t) ,
  \label{eq:ess_senoidal}
\end{equation}
expresión que muestra que el error es pequeño cuando $T(j\omega_0)$ es próximo a la unidad. Como se verá, con atascamiento-deslizamiento el error real lo domina el ciclo límite y no esta componente lineal.

\paragraph{Transiciones atascamiento-deslizamiento.}
Cada vez que la velocidad del levitante se anula, el sistema transita del régimen de deslizamiento al de atascamiento. Como se verá en los resultados, estas transiciones ocurren de forma continua, también lejos de los instantes en que la referencia invierte su sentido ($t = T/4,\,3T/4,\ldots$), y dan lugar a un ciclo límite.
El modelo de Karnopp detecta el atascamiento cuando se satisface la condición

\begin{equation}
  |v| \leq DV = 0{,}02 \text{ cm/s} ,
  \label{eq:banda_karnopp}
\end{equation}
y la fuerza magnética neta no supera la fricción estática máxima $F_s = 20$~kg\,cm/s$^2$.
En esa condición, tanto la velocidad como la aceleración se anulan y el error de seguimiento
crece transitoriamente.
Una vez que la fuerza de control acumulada supera $F_s$ (\emph{breakaway}), el sistema
abandona el estado de reposo y retoma el seguimiento de la referencia.
La acción integral del controlador es la que, durante cada atascamiento, modifica el voltaje hasta que la fuerza neta supera la fricción estática; la comparación con un controlador PI se presenta al final de esta sección.

\begin{figure}[ht]
    \centering
    % Estilo común de las figuras de seguimiento (pgfplots).
\pgfplotsset{
  seguimiento/.style={
    width=0.92\textwidth, height=5.4cm,
    xlabel={tiempo [\si{s}]},
    grid=major, grid style={gray!25},
    tick label style={font=\small}, label style={font=\small}, legend style={font=\small},
    /pgf/number format/use comma,
    scaled ticks=false,
    y tick label style={/pgf/number format/fixed},
    x tick label style={/pgf/number format/fixed, /pgf/number format/1000 sep={}},
    table/col sep=comma,
    every axis plot/.append style={line join=round},
  },
}
\definecolor{tray}{RGB}{31,94,166}
\definecolor{refc}{RGB}{40,40,40}
\definecolor{ctl2}{RGB}{196,96,40}

\begin{tikzpicture}
\begin{axis}[seguimiento, name=principal, xmin=0, xmax=2000, ymin=-3.15, ymax=-1.85,
             xlabel={}, ylabel={posición [\si{cm}]},
             legend style={at={(0.99,0.03)}, anchor=south east}, legend cell align=left]
  \addplot[tray, line width=1pt] table[x=t, y=x] {images/seguimiento_results/tikz/sen_posicion.csv};
  \addlegendentry{posición $x_1$}
  \addplot[refc, densely dashed, line width=0.8pt] table[x=t, y=r] {images/seguimiento_results/tikz/sen_posicion.csv};
  \addlegendentry{referencia $r$}
  \fill[ctl2, opacity=0.15] (axis cs:244,-3.15) rectangle (axis cs:256,-1.85);
\end{axis}
\begin{axis}[seguimiento, at={(principal.south west)}, anchor=north west, yshift=-0.45cm,
             height=4.2cm, xmin=244, xmax=256,
             ylabel={detalle [\si{cm}]}, axis background/.style={fill=ctl2!4}]
  \addplot[tray, line width=1pt] table[x=t, y=x] {images/seguimiento_results/tikz/sen_detalle.csv};
  \addplot[refc, densely dashed, line width=0.8pt] table[x=t, y=r] {images/seguimiento_results/tikz/sen_detalle.csv};
\end{axis}
\end{tikzpicture}

    \caption[Seguimiento senoidal: posición]{Seguimiento de la referencia senoidal. Arriba, posición del levitante (línea continua) y referencia (discontinua) durante dos periodos; la banda sombreada marca la ventana de detalle de abajo (244--256~s, máximo de la referencia), donde se observa el ciclo límite de atascamiento-deslizamiento alrededor de la referencia.}
    \label{fig:sen_pos}
\end{figure}

La figura~\ref{fig:sen_pos} muestra que, a la escala de la referencia, la posición sigue a la senoide sin desfase ni atenuación apreciables, en concordancia con $|T(j\omega_0)|\approx 1$. La ventana de detalle revela, sin embargo, que el levitante avanza en saltos y permanece atascado entre ellos, formando un ciclo límite de aproximadamente $0{,}02$~cm pico a pico y periodo cercano a 1~s.

\FloatBarrier

\begin{figure}[ht]
    \centering
    % Estilo común de las figuras de seguimiento (pgfplots).
\pgfplotsset{
  seguimiento/.style={
    width=0.92\textwidth, height=5.4cm,
    xlabel={tiempo [\si{s}]},
    grid=major, grid style={gray!25},
    tick label style={font=\small}, label style={font=\small}, legend style={font=\small},
    /pgf/number format/use comma,
    scaled ticks=false,
    y tick label style={/pgf/number format/fixed},
    x tick label style={/pgf/number format/fixed, /pgf/number format/1000 sep={}},
    table/col sep=comma,
    every axis plot/.append style={line join=round},
  },
}
\definecolor{tray}{RGB}{31,94,166}
\definecolor{refc}{RGB}{40,40,40}
\definecolor{ctl2}{RGB}{196,96,40}

\begin{tikzpicture}
\begin{axis}[seguimiento, name=principal, xmin=0, xmax=2000, xlabel={}, ylabel={velocidad [\si{cm/s}]}]
  \addplot[tray, line width=0.6pt] table[x=t, y=v] {images/seguimiento_results/tikz/sen_velocidad.csv};
  \fill[ctl2, opacity=0.15] ({axis cs:244,0} |- {rel axis cs:0,0}) rectangle ({axis cs:256,0} |- {rel axis cs:0,1});
\end{axis}
\begin{axis}[seguimiento, at={(principal.south west)}, anchor=north west, yshift=-0.45cm,
             height=4.2cm, xmin=244, xmax=256, ylabel={detalle [\si{cm/s}]}, axis background/.style={fill=ctl2!4}]
  \addplot[tray, line width=0.8pt] table[x=t, y=v] {images/seguimiento_results/tikz/sen_detalle.csv};
\end{axis}
\end{tikzpicture}

    \caption[Seguimiento senoidal: velocidad]{Velocidad del levitante durante el seguimiento senoidal. Los pulsos corresponden a los eventos de deslizamiento y los tramos en cero al régimen de atascamiento (75\,\% del tiempo).}
    \label{fig:sen_vel}
\end{figure}

La figura~\ref{fig:sen_vel} muestra que la velocidad alterna pulsos de deslizamiento de hasta $\pm 1$~cm/s con tramos de velocidad nula. En los 2000~s simulados se registraron 4072 rupturas (unas dos por segundo) y el levitante permaneció atascado el 75\,\% del tiempo. El atascamiento ocurre, por tanto, a lo largo de toda la prueba.

\FloatBarrier

\begin{figure}[ht]
    \centering
    % Estilo común de las figuras de seguimiento (pgfplots).
\pgfplotsset{
  seguimiento/.style={
    width=0.92\textwidth, height=5.4cm,
    xlabel={tiempo [\si{s}]},
    grid=major, grid style={gray!25},
    tick label style={font=\small}, label style={font=\small}, legend style={font=\small},
    /pgf/number format/use comma,
    scaled ticks=false,
    y tick label style={/pgf/number format/fixed},
    x tick label style={/pgf/number format/fixed, /pgf/number format/1000 sep={}},
    table/col sep=comma,
    every axis plot/.append style={line join=round},
  },
}
\definecolor{tray}{RGB}{31,94,166}
\definecolor{refc}{RGB}{40,40,40}
\definecolor{ctl2}{RGB}{196,96,40}

\begin{tikzpicture}
\begin{axis}[seguimiento, xmin=0, xmax=2000, ylabel={error $e=r-x_1$ [\si{cm}]}]
  \addplot[refc!60, thin] coordinates {(0,0) (2000,0)};
  \addplot[tray, line width=0.7pt] table[x=t, y=e] {images/seguimiento_results/tikz/sen_error.csv};
\end{axis}
\end{tikzpicture}

    \caption[Seguimiento senoidal: error]{Error de seguimiento $e(t)=r(t)-x_1(t)$ para la referencia senoidal.}
    \label{fig:sen_error}
\end{figure}

La figura~\ref{fig:sen_error} muestra que el error permanece acotado en una banda de $\pm 0{,}02$~cm (error máximo $0{,}0198$~cm, valor eficaz $0{,}0124$~cm) con media prácticamente nula. La amplitud de la banda varía ligeramente con la posición de la referencia, pero no se observan picos aislados: el error está dominado por el ciclo límite y no por la dinámica lineal de seguimiento.

\FloatBarrier

\begin{figure}[ht]
    \centering
    % Estilo común de las figuras de seguimiento (pgfplots).
\pgfplotsset{
  seguimiento/.style={
    width=0.92\textwidth, height=5.4cm,
    xlabel={tiempo [\si{s}]},
    grid=major, grid style={gray!25},
    tick label style={font=\small}, label style={font=\small}, legend style={font=\small},
    /pgf/number format/use comma,
    scaled ticks=false,
    y tick label style={/pgf/number format/fixed},
    x tick label style={/pgf/number format/fixed, /pgf/number format/1000 sep={}},
    table/col sep=comma,
    every axis plot/.append style={line join=round},
  },
}
\definecolor{tray}{RGB}{31,94,166}
\definecolor{refc}{RGB}{40,40,40}
\definecolor{ctl2}{RGB}{196,96,40}

\begin{tikzpicture}
\begin{axis}[seguimiento, name=principal, xmin=0, xmax=2000, ymin=-0.2, ymax=3.8, xlabel={},
             ylabel={señal de control [\si{V}]},
             legend style={at={(0.99,0.03)}, anchor=south east}, legend cell align=left, legend columns=3]
  \addplot[ctl2!45, line width=1.2pt] table[x=t, y=u_antes] {images/seguimiento_results/tikz/sen_control.csv};
  \addlegendentry{antes de la ZM}
  \addplot[tray, line width=0.5pt] table[x=t, y=u_desp] {images/seguimiento_results/tikz/sen_control.csv};
  \addlegendentry{después de la ZM}
  \addplot[red!70!black, densely dotted, line width=0.8pt] coordinates {(0,3.5) (2000,3.5)};
  \addlegendentry{$u_{\max}$}
  \fill[ctl2, opacity=0.15] (axis cs:244,-0.2) rectangle (axis cs:256,3.8);
\end{axis}
\begin{axis}[seguimiento, at={(principal.south west)}, anchor=north west, yshift=-0.45cm,
             height=4.2cm, xmin=244, xmax=256, ylabel={detalle [\si{V}]}, axis background/.style={fill=ctl2!4}]
  \addplot[tray, line width=0.8pt] table[x=t, y=u] {images/seguimiento_results/tikz/sen_detalle.csv};
\end{axis}
\end{tikzpicture}

    \caption[Seguimiento senoidal: señal de control]{Voltaje de control antes y después de la zona muerta durante el seguimiento senoidal. Abajo, detalle de 244 a 256~s: diente de sierra producido por la acción integral durante cada atascamiento.}
    \label{fig:sen_ctrl}
\end{figure}

La figura~\ref{fig:sen_ctrl} muestra que el voltaje sigue la forma de la referencia entre $1{,}33$ y $2{,}78$~V, siempre dentro del rango admisible $u \in [0;\,3{,}5]$~V y sin saturar. El detalle revela un perfil en diente de sierra: durante cada atascamiento la acción integral del PII hace crecer el voltaje hasta vencer la fricción estática, y tras la ruptura el voltaje cae bruscamente.

\FloatBarrier

%=============================================================================
\subsection{Seguimiento de referencia trapezoidal}
%=============================================================================

La referencia trapezoidal consiste en segmentos planos alternados en $x_1 = -2$ cm y
$x_1 = -3$ cm, conectados por rampas lineales de duración $250$ s cada una.
La señal describe así un perfil periódico de posición que combina fases de regulación
(segmentos planos) con fases de transición (rampas), constituyendo un escenario más exigente
que la referencia senoidal debido a la discontinuidad en la derivada de la referencia en los
instantes de inicio y fin de cada rampa.

\paragraph{Comportamiento durante la rampa.}
En la fase de rampa, la referencia varía a velocidad constante
$\dot{r} = (x_{f}-x_{i})/250 \text{ cm/s}$.
Según el modelo lineal, el error de seguimiento en esta fase converge a un valor constante determinado por el tipo del lazo.
En teoría, la segunda acción integral elimina el error de posición constante que un controlador de tipo uno mantendría durante la rampa; para la pendiente empleada, sin embargo, ese error sería del orden de $10^{-5}$~cm (véase la comparación al final de esta sección).

\paragraph{Comportamiento en los segmentos planos.}
Una vez que el levitante alcanza el nivel de referencia ($x_1 = -2$ cm o $x_1 = -3$ cm),
el sistema opera en modo de regulación.
Sin atascamiento-deslizamiento, la acción integral llevaría el error a cero; con él, como en la prueba de regulación, el error medio se anula pero el levitante oscila alrededor de la referencia en un ciclo límite.

\paragraph{Efecto de ruptura al inicio de cada rampa.}
El efecto de ruptura (\emph{breakaway}) se manifiesta al inicio de cada
rampa, cuando el sistema debe abandonar el estado de atascamiento en el que se encontraba
durante el segmento plano previo.
En ese instante, la fuerza magnética acumulada por la acción integral debe superar la
fricción estática máxima $F_s = 20$~kg\,cm/s$^2$ antes de que el levitante comience a desplazarse.
En la simulación, este retraso es de aproximadamente un segundo al inicio de la primera rampa y no produce un pico de error mayor que la banda del ciclo límite.

\paragraph{Señal de control durante la transición.}
En cada atascamiento, la acción integral del controlador modifica el voltaje hasta que la fuerza neta supera la fricción estática; tras la ruptura, el voltaje cae bruscamente, lo que produce un perfil en diente de sierra.
En todos los casos el voltaje se mantiene dentro del rango saturado $u \in [0;\,3{,}5]$ V,
sin alcanzar el límite superior.

\begin{figure}[ht]
    \centering
    % Estilo común de las figuras de seguimiento (pgfplots).
\pgfplotsset{
  seguimiento/.style={
    width=0.92\textwidth, height=5.4cm,
    xlabel={tiempo [\si{s}]},
    grid=major, grid style={gray!25},
    tick label style={font=\small}, label style={font=\small}, legend style={font=\small},
    /pgf/number format/use comma,
    scaled ticks=false,
    y tick label style={/pgf/number format/fixed},
    x tick label style={/pgf/number format/fixed, /pgf/number format/1000 sep={}},
    table/col sep=comma,
    every axis plot/.append style={line join=round},
  },
}
\definecolor{tray}{RGB}{31,94,166}
\definecolor{refc}{RGB}{40,40,40}
\definecolor{ctl2}{RGB}{196,96,40}

\begin{tikzpicture}
\begin{axis}[seguimiento, name=principal, xmin=0, xmax=2500, ymin=-3.15, ymax=-1.85,
             xlabel={}, ylabel={posición [\si{cm}]},
             legend style={at={(0.99,0.03)}, anchor=south east}, legend cell align=left]
  \addplot[tray, line width=1pt] table[x=t, y=x] {images/seguimiento_results/tikz/trap_posicion.csv};
  \addlegendentry{posición $x_1$}
  \addplot[refc, densely dashed, line width=0.8pt] table[x=t, y=r] {images/seguimiento_results/tikz/trap_posicion.csv};
  \addlegendentry{referencia $r$}
  \fill[ctl2, opacity=0.15] (axis cs:246,-3.15) rectangle (axis cs:266,-1.85);
\end{axis}
\begin{axis}[seguimiento, at={(principal.south west)}, anchor=north west, yshift=-0.45cm,
             height=4.2cm, xmin=246, xmax=266,
             ylabel={detalle [\si{cm}]}, axis background/.style={fill=ctl2!4}]
  \addplot[tray, line width=1pt] table[x=t, y=x] {images/seguimiento_results/tikz/trap_detalle.csv};
  \addplot[refc, densely dashed, line width=0.8pt] table[x=t, y=r] {images/seguimiento_results/tikz/trap_detalle.csv};
\end{axis}
\end{tikzpicture}

    \caption[Seguimiento trapezoidal: posición]{Seguimiento de la referencia trapezoidal. Arriba, posición y referencia durante dos periodos; abajo, detalle de 246 a 266~s, que incluye el inicio de la primera rampa en $t=250$~s.}
    \label{fig:trap_pos}
\end{figure}

La figura~\ref{fig:trap_pos} muestra que el levitante sigue las rampas y los segmentos planos sin error apreciable a la escala de la referencia. El detalle del inicio de la primera rampa muestra que el levitante permanece atascado alrededor de un segundo después de que la referencia comienza a moverse, y a partir de ahí sigue la rampa en escalones sucesivos de atascamiento-deslizamiento.

\FloatBarrier

\begin{figure}[ht]
    \centering
    % Estilo común de las figuras de seguimiento (pgfplots).
\pgfplotsset{
  seguimiento/.style={
    width=0.92\textwidth, height=5.4cm,
    xlabel={tiempo [\si{s}]},
    grid=major, grid style={gray!25},
    tick label style={font=\small}, label style={font=\small}, legend style={font=\small},
    /pgf/number format/use comma,
    scaled ticks=false,
    y tick label style={/pgf/number format/fixed},
    x tick label style={/pgf/number format/fixed, /pgf/number format/1000 sep={}},
    table/col sep=comma,
    every axis plot/.append style={line join=round},
  },
}
\definecolor{tray}{RGB}{31,94,166}
\definecolor{refc}{RGB}{40,40,40}
\definecolor{ctl2}{RGB}{196,96,40}

\begin{tikzpicture}
\begin{axis}[seguimiento, name=principal, xmin=0, xmax=2500, xlabel={}, ylabel={velocidad [\si{cm/s}]}]
  \addplot[tray, line width=0.6pt] table[x=t, y=v] {images/seguimiento_results/tikz/trap_velocidad.csv};
  \fill[ctl2, opacity=0.15] ({axis cs:246,0} |- {rel axis cs:0,0}) rectangle ({axis cs:266,0} |- {rel axis cs:0,1});
\end{axis}
\begin{axis}[seguimiento, at={(principal.south west)}, anchor=north west, yshift=-0.45cm,
             height=4.2cm, xmin=246, xmax=266, ylabel={detalle [\si{cm/s}]}, axis background/.style={fill=ctl2!4}]
  \addplot[tray, line width=0.8pt] table[x=t, y=v] {images/seguimiento_results/tikz/trap_detalle.csv};
\end{axis}
\end{tikzpicture}

    \caption[Seguimiento trapezoidal: velocidad]{Velocidad del levitante durante el seguimiento trapezoidal. El levitante permanece en reposo durante el primer segmento plano; a partir de la primera rampa, los pulsos de deslizamiento se mantienen durante el resto de la prueba.}
    \label{fig:trap_vel}
\end{figure}

La figura~\ref{fig:trap_vel} muestra que el levitante permanece en reposo durante el primer segmento plano, en el que parte del equilibrio. A partir de la primera rampa aparecen pulsos de deslizamiento de hasta $\pm 1$~cm/s, que continúan también en los segmentos planos posteriores: una vez iniciado, el ciclo límite no se extingue. En total se registraron 4400 rupturas en los 2500~s (dos periodos) y el levitante estuvo atascado el 78\,\% del tiempo.

\FloatBarrier

\begin{figure}[ht]
    \centering
    % Estilo común de las figuras de seguimiento (pgfplots).
\pgfplotsset{
  seguimiento/.style={
    width=0.92\textwidth, height=5.4cm,
    xlabel={tiempo [\si{s}]},
    grid=major, grid style={gray!25},
    tick label style={font=\small}, label style={font=\small}, legend style={font=\small},
    /pgf/number format/use comma,
    scaled ticks=false,
    y tick label style={/pgf/number format/fixed},
    x tick label style={/pgf/number format/fixed, /pgf/number format/1000 sep={}},
    table/col sep=comma,
    every axis plot/.append style={line join=round},
  },
}
\definecolor{tray}{RGB}{31,94,166}
\definecolor{refc}{RGB}{40,40,40}
\definecolor{ctl2}{RGB}{196,96,40}

\begin{tikzpicture}
\begin{axis}[seguimiento, xmin=0, xmax=2500, ylabel={error $e=r-x_1$ [\si{cm}]}]
  \addplot[refc!60, thin] coordinates {(0,0) (2500,0)};
  \addplot[tray, line width=0.7pt] table[x=t, y=e] {images/seguimiento_results/tikz/trap_error.csv};
\end{axis}
\end{tikzpicture}

    \caption[Seguimiento trapezoidal: error]{Error de seguimiento $e(t)=r(t)-x_1(t)$ para la referencia trapezoidal.}
    \label{fig:trap_error}
\end{figure}

La figura~\ref{fig:trap_error} muestra un error acotado en $\pm 0{,}022$~cm (valor eficaz $0{,}0114$~cm) y media nula, tanto en las rampas como en los segmentos planos. Al inicio de las rampas el error conserva la banda de amplitud casi uniforme que impone el ciclo límite, sin picos aislados.

\FloatBarrier

\begin{figure}[ht]
    \centering
    % Estilo común de las figuras de seguimiento (pgfplots).
\pgfplotsset{
  seguimiento/.style={
    width=0.92\textwidth, height=5.4cm,
    xlabel={tiempo [\si{s}]},
    grid=major, grid style={gray!25},
    tick label style={font=\small}, label style={font=\small}, legend style={font=\small},
    /pgf/number format/use comma,
    scaled ticks=false,
    y tick label style={/pgf/number format/fixed},
    x tick label style={/pgf/number format/fixed, /pgf/number format/1000 sep={}},
    table/col sep=comma,
    every axis plot/.append style={line join=round},
  },
}
\definecolor{tray}{RGB}{31,94,166}
\definecolor{refc}{RGB}{40,40,40}
\definecolor{ctl2}{RGB}{196,96,40}

\begin{tikzpicture}
\begin{axis}[seguimiento, name=principal, xmin=0, xmax=2500, ymin=-0.2, ymax=3.8, xlabel={},
             ylabel={señal de control [\si{V}]},
             legend style={at={(0.99,0.03)}, anchor=south east}, legend cell align=left, legend columns=3]
  \addplot[ctl2!45, line width=1.2pt] table[x=t, y=u_antes] {images/seguimiento_results/tikz/trap_control.csv};
  \addlegendentry{antes de la ZM}
  \addplot[tray, line width=0.5pt] table[x=t, y=u_desp] {images/seguimiento_results/tikz/trap_control.csv};
  \addlegendentry{después de la ZM}
  \addplot[red!70!black, densely dotted, line width=0.8pt] coordinates {(0,3.5) (2500,3.5)};
  \addlegendentry{$u_{\max}$}
  \fill[ctl2, opacity=0.15] (axis cs:246,-0.2) rectangle (axis cs:266,3.8);
\end{axis}
\begin{axis}[seguimiento, at={(principal.south west)}, anchor=north west, yshift=-0.45cm,
             height=4.2cm, xmin=246, xmax=266, ylabel={detalle [\si{V}]}, axis background/.style={fill=ctl2!4}]
  \addplot[tray, line width=0.8pt] table[x=t, y=u] {images/seguimiento_results/tikz/trap_detalle.csv};
\end{axis}
\end{tikzpicture}

    \caption[Seguimiento trapezoidal: señal de control]{Voltaje de control antes y después de la zona muerta durante el seguimiento trapezoidal. Abajo, detalle de 246 a 266~s.}
    \label{fig:trap_ctrl}
\end{figure}

La figura~\ref{fig:trap_ctrl} muestra que el voltaje sigue el perfil trapezoidal entre $1{,}32$ y $2{,}78$~V, sin saturar. El detalle muestra cómo, al inicio de la rampa, el voltaje desciende de forma suave hasta que la fuerza neta vence la fricción estática, y a partir de ese instante adopta el mismo perfil en diente de sierra observado en la prueba senoidal.

\FloatBarrier

%=============================================================================
\subsection{Análisis comparativo: PI frente a PII}
%=============================================================================

Las pruebas de regulación y seguimiento muestran que el PII estabiliza el sistema, pero no permiten saber qué aporta la segunda acción integral respecto a un controlador con una sola.

Para contrastar el papel de la doble acción integral, ambas pruebas se repitieron con el controlador PI, con una sola acción integral, $C_{PI}(s)=178{,}14\,(s+15)(s+12{,}91)/[s\,(s+400)]$, en las mismas condiciones: escalón de $-2$ a $-3$~cm y un periodo de la referencia trapezoidal, con zona muerta, fricción de Karnopp y $u\in[0;\,3{,}5]$~V. La tabla~\ref{tab:pi_pii} resume los resultados.

\begin{table}[ht]
\centering
\caption{Comparación entre los controladores PI y PII con atascamiento-deslizamiento.}
\label{tab:pi_pii}
\begin{tabular}{|l|c|c|}
\hline
 & PI & PII \\ \hline
Sobrepaso ante el escalón & $25{,}9$\,\% & $23{,}6$\,\% \\ \hline
Ciclo límite pico a pico (escalón) & $0{,}031$~cm & $0{,}033$~cm \\ \hline
Duración media de cada atascamiento (escalón) & $0{,}26$~s & $0{,}36$~s \\ \hline
Duración media de cada atascamiento (trapezoidal) & $0{,}29$~s & $0{,}39$~s \\ \hline
Atascamiento más largo (trapezoidal) & $1{,}07$~s & $3{,}13$~s \\ \hline
Fracción del tiempo atascado (trapezoidal) & $74$\,\% & $76$\,\% \\ \hline
Error eficaz en las rampas & $0{,}0112$~cm & $0{,}0118$~cm \\ \hline
Error medio en las rampas & $-3\times10^{-5}$~cm & $1\times10^{-5}$~cm \\ \hline
\end{tabular}
\end{table}

Los resultados no muestran la ventaja que se esperaba de la segunda acción integral. El PII reduce ligeramente el sobrepaso, pero los atascamientos son en promedio más largos que con el PI y el ciclo límite tiene prácticamente la misma amplitud. En las rampas, el error residual que en teoría produce un controlador de tipo uno es $\dot r/K_v \approx 4\times10^{-5}$~cm para la pendiente empleada ($0{,}004$~cm/s), con $K_v=91{,}6$~s$^{-1}$ (sección~\ref{sec:analisis_ciclo}), muy inferior a la amplitud del ciclo límite, por lo que con esta referencia no es posible distinguir entre ambos controladores. Ambos operan sin saturar y con errores eficaces del mismo orden.

Para descartar que esta diferencia se deba a la sintonización particular de cada controlador, se evaluaron dos familias de controladores construidas a partir de los anteriores: en cada una se escaló la ganancia por un factor entre $0{,}5$ y $3$ y se desplazaron los ceros de baja frecuencia (los asociados a la acción integral) por un factor entre $0{,}25$ y $4$. De las combinaciones resultantes, 256 son estables con un margen de fase de al menos $30^\circ$ en todo el intervalo de $-3$ a $-2$~cm; de ellas, 54 PI y 51 PII completan, sin saturar ni perder la estabilidad, el escalón de $-2$ a $-3$~cm y una rampa de $-0{,}05$~cm/s, doce veces más pronunciada que la de la referencia trapezoidal. La figura~\ref{fig:r15_familias} muestra, para cada uno, la amplitud del ciclo límite y la duración media de los atascamientos.

\begin{figure}[ht]
    \centering
    % Generado a partir de calculus/Final_Bien/R15_sintonizacion (familias_R15.m, simular_R15.py).
\definecolor{tray}{RGB}{31,94,166}
\definecolor{ctl2}{RGB}{196,96,40}
\begin{tikzpicture}
\begin{axis}[width=0.85\textwidth, height=7cm, xmin=0, xmax=0.12, ymin=0, ymax=1.5,
    xlabel={ciclo límite pico a pico [\si{cm}]}, ylabel={duración media del atascamiento [\si{s}]},
    grid=major, grid style={gray!25}, tick label style={font=\small}, label style={font=\small},
    legend style={font=\small, at={(0.98,0.98)}, anchor=north east}, legend cell align=left,
    /pgf/number format/use comma, scaled ticks=false, x tick label style={/pgf/number format/fixed, /pgf/number format/precision=2},
    table/col sep=comma]
  \addplot[only marks, mark=o, mark size=2.2pt, ctl2] table[x=ciclo, y=atasc] {images/r15_results/tikz/familia_pi.csv};
  \addlegendentry{familia PI}
  \addplot[only marks, mark=square, mark size=2pt, tray] table[x=ciclo, y=atasc] {images/r15_results/tikz/familia_pii.csv};
  \addlegendentry{familia PII}
  \addplot[only marks, mark=*, mark size=3.5pt, ctl2, draw=black] coordinates {(0.0309,0.259)};
  \addlegendentry{PI de la comparación}
  \addplot[only marks, mark=square*, mark size=3.2pt, tray, draw=black] coordinates {(0.0339,0.352)};
  \addlegendentry{PII de la tesis}
\end{axis}
\end{tikzpicture}

    \caption[Familias PI y PII]{Ciclo límite y duración media de los atascamientos en el escalón de $-2$ a $-3$~cm para las familias de controladores PI y PII que operan sin saturar ($u\in[0;\,3{,}5]$~V). Los símbolos rellenos corresponden a los controladores de la tabla~\ref{tab:pi_pii}.}
    \label{fig:r15_familias}
\end{figure}

El ciclo límite más pequeño de la familia PII ($0{,}026$~cm) duplica el del mejor PI ($0{,}013$~cm), y para amplitudes de ciclo límite similares los PI se atascan durante menos tiempo. La amplitud del ciclo límite disminuye al aumentar la ganancia proporcional, no al añadir una segunda integración. La ventaja de la doble acción integral sí aparece en la rampa: con los controladores de la tabla~\ref{tab:pi_pii}, el error medio es de $0{,}13\times10^{-3}$~cm con el PII y de $0{,}56\times10^{-3}$~cm con el PI. Sin embargo, ambos son más de diez veces menores que el error eficaz que impone el ciclo límite en esa prueba, de unos $7{,}6\times10^{-3}$~cm y prácticamente igual con los dos controladores, por lo que no se traducen en un mejor seguimiento. La validación con integración \texttt{ode45} de los controladores extremos reproduce estos valores con diferencias menores al 8\,\%.

En síntesis, los experimentos de regulación y seguimiento demuestran que el controlador PII
diseñado es capaz de estabilizar el sistema de levitación magnética en configuración inestable
y de mantener un desempeño de seguimiento aceptable bajo los efectos combinados de la
zona muerta dinámica del actuador y la fricción de Karnopp, sin necesidad de compensación
explícita de ninguna de estas dos no linealidades.

% Análisis del ciclo límite. Datos: calculus/Final_Bien/R15_sintonizacion/
% (analisis_ciclo.m, simular_R15.py, exportar_ciclo.py).
\section{Análisis del ciclo límite}
\label{sec:analisis_ciclo}

Las pruebas anteriores muestran que, en presencia de atascamiento-deslizamiento, el levitante nunca queda en reposo sobre la referencia: oscila alrededor de ella con una amplitud de unas centésimas de centímetro. También muestran que la segunda acción integral del PII no reduce esa oscilación respecto a un PI. Esta sección explica ambos hechos paso a paso. Primero se traduce la fricción estática a términos de voltaje; después se describe un ciclo completo, se identifica la acción integral como su causa y se estiman la duración de cada atascamiento y la amplitud del ciclo. Al final se discute el único escenario en el que la segunda integración sí ofrece una ventaja. Todas las cifras corresponden al escalón de $-2$ a $-3$~cm con atascamiento-deslizamiento; las estadísticas de los ciclos se calculan en el estado estacionario, entre 20 y 40~s, por lo que difieren ligeramente de las de la tabla~\ref{tab:pi_pii}, que abarcan todo el intervalo posterior al escalón.

%-----------------------------------------------------------------------------
\subsection{La fricción estática vista desde el voltaje}

Mientras el levitante está atascado, el modelo de Karnopp lo mantiene inmóvil siempre que la fuerza neta $F_e = u/[b(a-x)^4] - mg$ no supere la fricción estática máxima $F_s$. Despejando el voltaje, el imán permanece atascado en la posición $x$ mientras
\begin{equation}
u_e(x)\left(1-\frac{F_s}{mg}\right) \;\le\; u \;\le\; u_e(x)\left(1+\frac{F_s}{mg}\right) .
\label{eq:banda_voltaje}
\end{equation}
El centro de la banda~\eqref{eq:banda_voltaje} es el voltaje de equilibrio $u_e(x)=mgb(a-x)^4$, el que sostiene al imán en la posición $x$. La fricción estática se manifiesta así como una banda muerta de voltaje centrada en $u_e$, con una anchura relativa de $\pm F_s/(mg)=\pm14{,}5$\,\%. En $x_1=-3$~cm, donde $u_e=2{,}393$~V, la banda va de $2{,}047$ a $2{,}739$~V: mientras el voltaje se mantenga dentro de ese intervalo de $0{,}69$~V, el imán no se mueve. Para despegarlo, el controlador tiene que llevar el voltaje hasta uno de los bordes, es decir, cambiarlo al menos $u_e F_s/(mg)=0{,}346$~V respecto al equilibrio.

%-----------------------------------------------------------------------------
\subsection{Anatomía de un ciclo}

La figura~\ref{fig:anatomia} muestra tres segundos del estado estacionario con el PII. Cada ciclo consta de cuatro fases que se repiten de forma casi idéntica.

\begin{enumerate}
  \item \textbf{Atascamiento.} El levitante queda detenido a un lado de la referencia, con un error de unos $0{,}016$~cm. La velocidad es nula y el error es constante, pero no nulo, de modo que la acción integral sigue modificando el voltaje: en la gráfica inferior se observa una rampa que se aleja de $u_e$.
  \item \textbf{Ruptura.} Cuando el voltaje alcanza el borde de la banda~\eqref{eq:banda_voltaje}, la fuerza neta supera $F_s$ y el levitante se despega. En promedio, el voltaje cambia $0{,}34$~V durante cada atascamiento, prácticamente la mitad de la anchura de la banda.
  \item \textbf{Deslizamiento.} Una vez en movimiento, la fricción pasa de estática a viscosa y desaparece de golpe la fuerza que equilibraba al imán. La fuerza en exceso lo acelera hacia la referencia, y el controlador no puede retirar instantáneamente el voltaje que acumuló durante el atascamiento: el levitante cruza la referencia y recorre unos $0{,}032$~cm.
  \item \textbf{Nuevo atascamiento.} Cuando la velocidad vuelve a anularse, el imán queda detenido del otro lado de la referencia y el proceso se repite en sentido contrario.
\end{enumerate}

Cada ciclo completo tiene dos atascamientos y dura en promedio $0{,}98$~s con el PII y $0{,}74$~s con el PI. La amplitud pico a pico ($0{,}033$ y $0{,}031$~cm, respectivamente) coincide con la distancia recorrida en cada deslizamiento.

\begin{figure}[ht]
    \centering
    % Estilo común de las figuras del análisis del ciclo límite (pgfplots).
\pgfplotsset{
  ciclo/.style={
    width=0.9\textwidth, height=4.2cm,
    grid=major, grid style={gray!25},
    tick label style={font=\small}, label style={font=\small}, legend style={font=\small},
    /pgf/number format/use comma, scaled ticks=false,
    y tick label style={/pgf/number format/fixed},
    x tick label style={/pgf/number format/fixed},
    table/col sep=comma, unbounded coords=jump,
    every axis plot/.append style={line join=round},
  },
}
\definecolor{tray}{RGB}{31,94,166}
\definecolor{refc}{RGB}{40,40,40}
\definecolor{ctl2}{RGB}{196,96,40}
\definecolor{ver}{RGB}{46,133,64}

\begin{tikzpicture}
\begin{axis}[ciclo, name=p1, xmin=20, xmax=23, xticklabels={}, ylabel={$x_1$ [\si{cm}]}, ymin=-3.025, ymax=-2.975,
    y tick label style={/pgf/number format/precision=2}]
  \addplot[refc, densely dashed] coordinates {(20,-3) (23,-3)};
  \addplot[tray, line width=1pt] table[x=t, y=x] {images/ciclo_results/tikz/anatomia_pii.csv};
\end{axis}
\begin{axis}[ciclo, name=p2, at={(p1.south west)}, anchor=north west, yshift=-0.25cm, xmin=20, xmax=23, xticklabels={},
    ylabel={$x_2$ [\si{cm/s}]}]
  \addplot[tray, line width=0.8pt] table[x=t, y=v] {images/ciclo_results/tikz/anatomia_pii.csv};
\end{axis}
\begin{axis}[ciclo, at={(p2.south west)}, anchor=north west, yshift=-0.25cm, xmin=20, xmax=23, ymin=1.88, ymax=2.92,
    xlabel={tiempo [\si{s}]}, ylabel={$u$ [\si{V}]}]
  \addplot[ver, densely dashed, line width=0.9pt] coordinates {(20,2.0473) (23,2.0473)};
  \addplot[ver, densely dashed, line width=0.9pt] coordinates {(20,2.7394) (23,2.7394)};
  \addplot[refc!60, densely dotted] coordinates {(20,2.3934) (23,2.3934)};
  \addplot[tray, line width=1pt] table[x=t, y=u] {images/ciclo_results/tikz/anatomia_pii.csv};
  \node[font=\scriptsize, anchor=south east, ver!70!black] at (axis cs:22.98,2.7394) {$u_e(1+F_s/mg)$};
  \node[font=\scriptsize, anchor=north east, ver!70!black] at (axis cs:22.98,2.0473) {$u_e(1-F_s/mg)$};
  \node[font=\scriptsize, anchor=south west, refc!80] at (axis cs:20.02,2.3934) {$u_e$};
\end{axis}
\end{tikzpicture}

    \caption[Anatomía del ciclo límite]{Tres segundos del ciclo límite con el PII en $x_1=-3$~cm. De arriba abajo: posición, velocidad y voltaje. Las líneas verdes discontinuas son los bordes de la banda~\eqref{eq:banda_voltaje}: el levitante solo se despega cuando el voltaje llega a uno de ellos.}
    \label{fig:anatomia}
\end{figure}

%-----------------------------------------------------------------------------
\subsection{La acción integral es la causa del ciclo}

El mecanismo anterior depende de que el voltaje siga cambiando mientras el levitante está detenido. En un controlador sin acción integral, la salida depende solo del error actual y de su derivada, y ambos permanecen constantes durante el atascamiento. Para comprobarlo, se suprimieron del PII los dos integradores y los dos ceros de baja frecuencia que los acompañan, lo que deja el controlador PD
\begin{equation}
C_{PD}(s) = \frac{152{,}43\,(s+17{,}31)}{s+301{,}1} ,
\label{eq:pd}
\end{equation}
que estabiliza el lazo en todo el intervalo de operación. Como un PD no puede generar por sí mismo el voltaje de equilibrio, se le sumó la prealimentación $u_e(r)$ correspondiente a la referencia.

La figura~\ref{fig:pd_pi_pii} compara los tres controladores. Con el PD, el levitante llega a la nueva referencia, se atasca y permanece inmóvil, sin ciclo límite. A cambio, el error final puede ser distinto de cero. La ganancia estática del PD, $8{,}76$~V/cm, equivale a una rigidez de $506{,}5$~kg\,cm/s$^2$ por centímetro de error, de modo que el levitante puede quedar detenido en cualquier punto donde el error sea menor que $F_s/506{,}5 = 0{,}039$~cm. En esta simulación se detuvo a $0{,}002$~cm de la referencia.

\begin{figure}[ht]
    \centering
    % Estilo común de las figuras del análisis del ciclo límite (pgfplots).
\pgfplotsset{
  ciclo/.style={
    width=0.9\textwidth, height=4.2cm,
    grid=major, grid style={gray!25},
    tick label style={font=\small}, label style={font=\small}, legend style={font=\small},
    /pgf/number format/use comma, scaled ticks=false,
    y tick label style={/pgf/number format/fixed},
    x tick label style={/pgf/number format/fixed},
    table/col sep=comma, unbounded coords=jump,
    every axis plot/.append style={line join=round},
  },
}
\definecolor{tray}{RGB}{31,94,166}
\definecolor{refc}{RGB}{40,40,40}
\definecolor{ctl2}{RGB}{196,96,40}
\definecolor{ver}{RGB}{46,133,64}

\begin{tikzpicture}
\begin{axis}[ciclo, name=p1, height=5cm, xmin=0, xmax=40, ymin=-3.3, ymax=-1.9, xlabel={},
    ylabel={posición [\si{cm}]}, legend style={at={(0.99,0.97)}, anchor=north east}, legend cell align=left, legend columns=3]
  \addplot[ver, line width=1.1pt] table[x=t, y=x_pd] {images/ciclo_results/tikz/pd_pi_pii.csv}; \addlegendentry{PD + $u_e(r)$}
  \addplot[ctl2, line width=0.7pt] table[x=t, y=x_pi] {images/ciclo_results/tikz/pd_pi_pii.csv}; \addlegendentry{PI}
  \addplot[tray, line width=0.7pt] table[x=t, y=x_pii] {images/ciclo_results/tikz/pd_pi_pii.csv}; \addlegendentry{PII}
  \fill[ctl2, opacity=0.12] (axis cs:15,-3.3) rectangle (axis cs:25,-1.9);
\end{axis}
\begin{axis}[ciclo, at={(p1.south west)}, anchor=north west, yshift=-0.45cm, height=4.2cm, xmin=15, xmax=25,
    ymin=-3.025, ymax=-2.975, xlabel={tiempo [\si{s}]}, ylabel={detalle [\si{cm}]}, y tick label style={/pgf/number format/precision=2}]
  \addplot[refc, densely dashed] coordinates {(15,-3) (25,-3)};
  \addplot[ctl2, line width=0.7pt] table[x=t, y=x_pi] {images/ciclo_results/tikz/pd_pi_pii.csv};
  \addplot[tray, line width=0.7pt] table[x=t, y=x_pii] {images/ciclo_results/tikz/pd_pi_pii.csv};
  \addplot[ver, line width=1.4pt] table[x=t, y=x_pd] {images/ciclo_results/tikz/pd_pi_pii.csv};
\end{axis}
\end{tikzpicture}

    \caption[PD, PI y PII con atascamiento-deslizamiento]{Escalón de $-2$ a $-3$~cm con el PD con prealimentación, el PI y el PII. Abajo, detalle de 15 a 25~s: el PD queda en reposo y los controladores con acción integral oscilan alrededor de la referencia.}
    \label{fig:pd_pi_pii}
\end{figure}

Este es el dilema clásico del control con fricción estática, conocido en la literatura como \emph{hunting}~\cite{olsson}. Sin acción integral, el lazo se detiene con un error que depende de dónde quedó atascado y de qué tan exacto sea el voltaje de equilibrio que se suministra. Con acción integral, el error medio se anula aunque el modelo no sea exacto, pero el lazo nunca se detiene: el integrador sigue acumulando el error residual hasta forzar una nueva ruptura. El control conmutado de Alcántara et al.~\cite{alcantara} aborda precisamente este dilema, desactivando la acción integral cuando el error es pequeño.

%-----------------------------------------------------------------------------
\subsection{Duración del atascamiento}

Si el error con el que el levitante se atasca es $e_0$, la duración del atascamiento es el tiempo que tarda la acción integral en cambiar el voltaje los $0{,}346$~V necesarios para la ruptura. A bajas frecuencias, el PI se comporta como $K_i/s$ con $K_i=86{,}2$~V/(cm\,s), y el PII como $K_{ii}/s^2 + K_i/s$ con $K_{ii}=79{,}5$~V/(cm\,s$^2$) y $K_i=63{,}4$~V/(cm\,s). Con el error constante, el cambio de voltaje durante el atascamiento es
\begin{equation}
\Delta u_{PI}(t) = K_i e_0\, t, \qquad
\Delta u_{PII}(t) = K_i e_0\, t + \frac{K_{ii} e_0}{2}\, t^2 .
\label{eq:rampa_atasco}
\end{equation}

Para el PI, con $e_0=0{,}0151$~cm, el voltaje crece a $K_i e_0 = 1{,}30$~V/s y alcanza el borde de la banda, que en promedio exige un cambio de $0{,}328$~V, en $0{,}328/1{,}30=0{,}25$~s, en acuerdo con la duración media medida de $0{,}26$~s. Para el PII, la figura~\ref{fig:atasco} muestra que el voltaje sigue la parábola del cambio de voltaje del PII~\eqref{eq:rampa_atasco}: la curvatura medida coincide con $K_{ii}e_0$ dentro de un 10\,\%. Sin embargo, el término cuadrático solo domina sobre el lineal cuando $t > 2K_i/K_{ii} = 1{,}6$~s, y los atascamientos duran unos $0{,}35$~s. En ese intervalo, la segunda integración aporta aproximadamente una quinta parte del cambio de voltaje, y el término lineal del PII es más lento que el del PI porque su $K_i$ es menor (63 frente a 86~V/(cm\,s)). Además, la pendiente inicial medida es un 20 a 25\,\% menor que $K_i e_0$, lo que indica que el estado del controlador al inicio de cada atascamiento conserva parte de la dinámica del deslizamiento previo. El resultado neto es que el PII tarda más que el PI en despegar el imán.

\begin{figure}[ht]
    \centering
    % Estilo común de las figuras del análisis del ciclo límite (pgfplots).
\pgfplotsset{
  ciclo/.style={
    width=0.9\textwidth, height=4.2cm,
    grid=major, grid style={gray!25},
    tick label style={font=\small}, label style={font=\small}, legend style={font=\small},
    /pgf/number format/use comma, scaled ticks=false,
    y tick label style={/pgf/number format/fixed},
    x tick label style={/pgf/number format/fixed},
    table/col sep=comma, unbounded coords=jump,
    every axis plot/.append style={line join=round},
  },
}
\definecolor{tray}{RGB}{31,94,166}
\definecolor{refc}{RGB}{40,40,40}
\definecolor{ctl2}{RGB}{196,96,40}
\definecolor{ver}{RGB}{46,133,64}

\begin{tikzpicture}
\begin{axis}[ciclo, height=6.2cm, xmin=0, xmax=0.45, ymin=0, ymax=0.42,
    xlabel={tiempo desde el inicio del atascamiento [\si{s}]}, ylabel={cambio de voltaje [\si{V}]},
    legend style={at={(0.98,0.03)}, anchor=south east}, legend cell align=left]
  \addplot[ver, densely dashed, line width=0.9pt] coordinates {(0,0.3461) (0.45,0.3461)};
  \addlegendentry{$u_e F_s/(mg)=0{,}346$ V}
  \addplot[ctl2, line width=0.8pt] table[x=tau, y=du] {images/ciclo_results/tikz/atasco_pi.csv};
  \addlegendentry{PI (simulación)}
  \addplot[ctl2!70!black, dashed, line width=1pt] table[x=tau, y=du] {images/ciclo_results/tikz/atasco_pi_teo.csv};
  \addlegendentry{PI: $K_i e_0 t$}
  \addplot[tray, line width=0.8pt] table[x=tau, y=du] {images/ciclo_results/tikz/atasco_pii.csv};
  \addlegendentry{PII (simulación)}
  \addplot[tray!70!black, dashed, line width=1pt] table[x=tau, y=du] {images/ciclo_results/tikz/atasco_pii_teo.csv};
  \addlegendentry{PII: $K_i e_0 t + K_{ii} e_0 t^2/2$}
\end{axis}
\end{tikzpicture}

    \caption[Voltaje durante el atascamiento]{Cambio de voltaje durante cada atascamiento, alineado al instante en que comienza y con el signo del error. Las curvas discontinuas son la predicción~\eqref{eq:rampa_atasco} con el error medio de atascamiento. La ruptura ocurre al alcanzar $u_eF_s/(mg)$.}
    \label{fig:atasco}
\end{figure}

%-----------------------------------------------------------------------------
\subsection{Amplitud del ciclo}

La amplitud del ciclo la determina la fase de deslizamiento. En el instante de la ruptura, el voltaje está en el borde de la banda, de modo que la fuerza neta sobre el imán vale aproximadamente $F_s$ y deja de estar compensada por la fricción. El imán se desplaza hasta que el controlador, a través de su respuesta a la posición, retira esa fuerza en exceso. Cuanto más rígido es el lazo, es decir, cuanto mayor es la ganancia que traduce el error de posición en voltaje, menor es el desplazamiento necesario. Este razonamiento predice que la amplitud del ciclo es aproximadamente inversamente proporcional a la ganancia del controlador, e independiente del número de integradores.

La figura~\ref{fig:ganancia} lo confirma con las familias de controladores de la figura~\ref{fig:r15_familias}, manteniendo fijos los ceros de baja frecuencia ($\alpha=1$) y variando solo el factor de ganancia $k_K$. La amplitud decrece aproximadamente como $k_K^{-1{,}2}$ en la familia PI y como $k_K^{-1{,}1}$ en la familia PII. Con la misma ganancia, ambas estructuras producen prácticamente el mismo ciclo límite: $0{,}072$ frente a $0{,}074$~cm con $k_K=0{,}5$, y $0{,}031$ frente a $0{,}034$~cm con $k_K=1$ (valores del barrido en la unidad de procesamiento gráfico, GPU, que difieren en menos de 3\,\% de los obtenidos con \texttt{ode45}). La ganancia fija la amplitud del ciclo; la segunda integración no la modifica.

\begin{figure}[ht]
    \centering
    % Estilo común de las figuras del análisis del ciclo límite (pgfplots).
\pgfplotsset{
  ciclo/.style={
    width=0.9\textwidth, height=4.2cm,
    grid=major, grid style={gray!25},
    tick label style={font=\small}, label style={font=\small}, legend style={font=\small},
    /pgf/number format/use comma, scaled ticks=false,
    y tick label style={/pgf/number format/fixed},
    x tick label style={/pgf/number format/fixed},
    table/col sep=comma, unbounded coords=jump,
    every axis plot/.append style={line join=round},
  },
}
\definecolor{tray}{RGB}{31,94,166}
\definecolor{refc}{RGB}{40,40,40}
\definecolor{ctl2}{RGB}{196,96,40}
\definecolor{ver}{RGB}{46,133,64}

\begin{tikzpicture}
\begin{loglogaxis}[ciclo, height=6cm, width=0.75\textwidth, xmin=0.45, xmax=1.1, ymin=0.025, ymax=0.085,
    xlabel={factor de ganancia $k_K$}, ylabel={ciclo límite pico a pico [\si{cm}]},
    xtick={0.5,0.6,0.7,0.8,0.9,1}, ytick={0.03,0.04,0.05,0.06,0.08}, log ticks with fixed point,
    legend style={at={(0.98,0.98)}, anchor=north east}, legend cell align=left]
  \addplot[ctl2, mark=o, mark size=2.2pt, line width=0.8pt] table[x=kK, y=ciclo] {images/ciclo_results/tikz/ganancia_pi.csv}; \addlegendentry{familia PI ($\alpha=1$)}
  \addplot[tray, mark=square, mark size=2pt, line width=0.8pt] table[x=kK, y=ciclo] {images/ciclo_results/tikz/ganancia_pii.csv}; \addlegendentry{familia PII ($\alpha=1$)}
  \addplot[refc!60, densely dashed, domain=0.45:1.1, samples=2] {0.032/x}; \addlegendentry{$\propto 1/k_K$}
\end{loglogaxis}
\end{tikzpicture}

    \caption[Ciclo límite contra la ganancia]{Amplitud del ciclo límite contra el factor de ganancia $k_K$, con los ceros de baja frecuencia fijos. Ganancias mayores saturan el actuador.}
    \label{fig:ganancia}
\end{figure}

%-----------------------------------------------------------------------------
\subsection{Dónde sí ayuda la segunda integración}

La ventaja teórica del PII aparece al seguir una rampa. El lazo con el PI es de tipo uno: ante una rampa de pendiente $\dot r$ conserva un error estacionario $\dot r/K_v$, donde $K_v=\lim_{s\to0} sC(s)G(s)=91{,}6$~s$^{-1}$ en $x_1^*=-3$~cm. El lazo con el PII es de tipo dos y anula ese error. Con la pendiente de la referencia trapezoidal ($0{,}004$~cm/s), el error del PI sería de solo $4\times10^{-5}$~cm. Con una rampa doce veces más pronunciada ($0{,}05$~cm/s), la predicción es $5{,}5\times10^{-4}$~cm, y la simulación da $5{,}6\times10^{-4}$~cm con el PI frente a $1{,}3\times10^{-4}$~cm con el PII. La segunda integración cumple su función, pero el ciclo límite impone en ambos casos un error eficaz de unos $7{,}6\times10^{-3}$~cm, más de diez veces mayor. En este sistema, la ventaja del PII solo sería apreciable con referencias de pendiente mucho mayor o con una fricción estática mucho menor.

%-----------------------------------------------------------------------------
\subsection{Síntesis}

\begin{itemize}
  \item La fricción estática equivale a una banda muerta de voltaje de $\pm14{,}5$\,\% alrededor de $u_e$; el imán solo se despega cuando el voltaje alcanza uno de sus bordes.
  \item La acción integral es la causa del ciclo límite: sin ella el levitante se detiene con un error acotado por $F_s$ y la rigidez del lazo; con ella el error medio se anula, pero el integrador provoca rupturas sucesivas.
  \item La duración de cada atascamiento depende del término integral simple, $K_i e_0$. El término doble del PII solo dominaría en atascamientos de más de $1{,}6$~s, y los observados duran unos $0{,}35$~s.
  \item La amplitud del ciclo es aproximadamente inversamente proporcional a la ganancia del controlador y no depende del número de integradores.
  \item La segunda integración elimina el error en rampa, pero ese error es más de diez veces menor que el que impone el ciclo límite, incluso con pendientes doce veces mayores que las de las pruebas de seguimiento.
\end{itemize}

\FloatBarrier

% Generado a partir de calculus/Final_Bien/evidencia_P8.m (simulador corregido).
\section{Configuración estable}

En la configuración repulsiva actúa la bobina inferior. Con $x$ medida hacia arriba, la distancia entre el imán y esa bobina es $a+x$, la fuerza electromagnética es $u/[b(a+x)^4]$ y el modelo linealizado tiene la matriz de estado $\bigl[\begin{smallmatrix}0&1\\-4g/(a+x_1^*)&-c_1\end{smallmatrix}\bigr]$, cuyos polos son complejos con parte real negativa: $-4{,}28\pm21{,}49j$ en $x_1^*=1$~cm y $-4{,}28\pm20{,}24j$ en $x_1^*=2$~cm (factor de amortiguamiento cercano a $0{,}2$). Se evaluaron en esta configuración tres controladores diseñados durante el trabajo: el PI de la comparación de la tabla~\ref{tab:pi_pii}, una versión preliminar del PII y el PII definitivo $C(s)$ de la sección~\ref{sec:pii}. Con la planta linealizada en $x_1^*=1{,}5$~cm, los tres estabilizan el lazo, con márgenes de fase de $50{,}7^\circ$, $51{,}0^\circ$ y $43{,}3^\circ$, respectivamente. Es decir, el mismo controlador diseñado para la configuración inestable estabiliza también la estable, aunque con un margen de fase menor.

\begin{table}[ht]
\centering
\caption{Configuración estable: escalón de 1 a 2~cm con zona muerta, fricción de Karnopp y $u\in[0;\,3{,}5]$~V.}
\label{tab:p8_estable}
\begin{tabular}{|l|c|c|c|}
\hline
 & PI & PII preliminar & PII (tesis) \\ \hline
Margen de fase (lineal, $x_1^*=1{,}5$~cm) & $50{,}7^\circ$ & $51{,}0^\circ$ & $43{,}3^\circ$ \\ \hline
Sobrepaso (lineal) & $19{,}7$\,\% & $18{,}6$\,\% & $26{,}8$\,\% \\ \hline
Sobrepaso (no lineal con AD) & $10{,}5$\,\% & $11{,}1$\,\% & $9{,}8$\,\% \\ \hline
Entrada en la banda de $\pm5$\,\% & $0{,}24$~s & $0{,}36$~s & $0{,}35$~s \\ \hline
Ciclo límite pico a pico & $0{,}018$~cm & $0{,}023$~cm & $0{,}020$~cm \\ \hline
\end{tabular}
\end{table}

La figura~\ref{fig:p8_estable} y la tabla~\ref{tab:p8_estable} muestran que los tres controladores llevan el levitante de 1 a 2~cm y que, como en la configuración atractiva, la zona muerta y la fricción estática producen un ciclo límite en estado estacionario, aquí de menor amplitud. En el transitorio la señal de control alcanza brevemente la saturación de $3{,}5$~V.

\begin{figure}[ht]
    \centering
    % Estilo común de las figuras de evidencia P8 (pgfplots).
\pgfplotsset{
  pocho/.style={
    width=0.92\textwidth, height=5.2cm,
    grid=major, grid style={gray!25},
    tick label style={font=\small}, label style={font=\small}, legend style={font=\small},
    /pgf/number format/use comma,
    scaled ticks=false,
    y tick label style={/pgf/number format/fixed},
    table/col sep=comma, unbounded coords=jump,
    every axis plot/.append style={line join=round},
  },
}
\definecolor{tray}{RGB}{31,94,166}
\definecolor{refc}{RGB}{40,40,40}
\definecolor{ctl2}{RGB}{196,96,40}
\definecolor{ver}{RGB}{46,133,64}

\begin{tikzpicture}
\begin{axis}[pocho, name=principal, xmin=0, xmax=40, ymin=0.8, ymax=2.3, xlabel={},
    ylabel={posición [\si{cm}]}, legend style={at={(0.99,0.03)}, anchor=south east}, legend cell align=left]
  \addplot[ctl2, line width=0.8pt] table[x=t, y=x_pi] {images/p8_results/tikz/estable.csv}; \addlegendentry{PI}
  \addplot[tray, line width=0.8pt] table[x=t, y=x_pii] {images/p8_results/tikz/estable.csv}; \addlegendentry{PII (tesis)}
\end{axis}
\begin{axis}[pocho, at={(principal.south west)}, anchor=north west, yshift=-0.45cm, height=4.4cm,
    xmin=0, xmax=40, ymin=-0.2, ymax=3.8, xlabel={tiempo [\si{s}]}, ylabel={$u$ [\si{V}]}]
  \addplot[ctl2, line width=0.7pt] table[x=t, y=u_pi] {images/p8_results/tikz/estable.csv};
  \addplot[tray, line width=0.7pt] table[x=t, y=u_pii] {images/p8_results/tikz/estable.csv};
  \addplot[red!70!black, densely dotted, line width=0.8pt] coordinates {(0,3.5) (40,3.5)};
\end{axis}
\end{tikzpicture}

    \caption[Configuración estable]{Regulación en la configuración estable (repulsiva) con el PI y con el PII de la tesis, escalón de 1 a 2~cm. Arriba, posición; abajo, señal de control.}
    \label{fig:p8_estable}
\end{figure}

\FloatBarrier

